% !TEX root = ../main.tex

\chapter{硬件平台参数}

\section{Unitree G1 机器人平台}

本文实验所涉及的人形机器人平台以 Unitree G1 为主。为便于后续补充正式实验配置，本节先给出参数表的组织方式，并保留关键字段占位。后续可根据实际平台规格填写自由度数量、机身尺寸、关节力矩范围、传感器配置和机载计算资源等信息。

\begin{table}[htbp]
  \centering
  \caption{Unitree G1 平台关键参数占位表}
  \begin{tabular}{ll}
    \toprule
    参数项 & 参数值 \\
    \midrule
    总自由度 & 待补充 \\
    机身高度/重量 & 待补充 \\
    下肢关键关节配置 & 待补充 \\
    上肢关键关节配置 & 待补充 \\
    IMU/编码器配置 & 待补充 \\
    机载计算单元 & 待补充 \\
    \bottomrule
  \end{tabular}
\end{table}

\section{遥操作与感知设备}

除机器人本体外，本文还涉及多种遥操作和感知设备，包括 PICO VR 设备、Noitom 惯性动捕系统、Nokov 光学动捕系统以及 RealSense D455 深度相机等。不同设备分别承担人体动作采集、环境感知和任务交互等功能。建议后续在正式定稿时补充其采样频率、视场角、定位精度、通信时延和标定方式等参数，以支撑实验复现。

\chapter{强化学习训练超参数}

\section{PPO 训练配置}

为便于后续复现实验，本文在本附录中预留强化学习训练超参数表，用于汇总第 3 章下肢 locomotion 策略所采用的核心配置。当前初稿阶段暂不填入未经确认的具体数值，仅给出表项结构。

\begin{table}[htbp]
  \centering
  \caption{PPO 训练超参数占位表}
  \begin{tabular}{ll}
    \toprule
    超参数项 & 参数值 \\
    \midrule
    学习率 & 待补充 \\
    折扣因子 $\gamma$ & 待补充 \\
    GAE 参数 $\lambda$ & 待补充 \\
    批次大小 & 待补充 \\
    每轮采样步数 & 待补充 \\
    裁剪系数 & 待补充 \\
    熵正则系数 & 待补充 \\
    \bottomrule
  \end{tabular}
\end{table}

\section{网络结构与训练环境}

除 PPO 核心超参数外，还需要对策略网络结构、价值网络结构、训练轮数、域随机化范围和奖励权重进行补充说明。后续可在本节中加入网络层数、隐藏层维度、激活函数类型、仿真时间步长以及各奖励项权重表格，以形成完整的训练配置说明。
